\scanpage{110}
\mspage{54}
l'ensemble fini
\[
\begin{aligned}
(154)\quad H(J)={}&\text{ensemble des plongements d'ens. ordonnés, compatibles avec}\\
&\text{les f codim, de }J\text{ dans }\deleted{\mathcal P([1,n]\cap\mathbb N)}\mathcal P(I_d)^\circ\\
&\text{(où }I_d\overset{\mathrm{def}}=[1,d]\cap\mathbb N\text{)}
\end{aligned}
\]
alors $\mathfrak S_d$ opère à gauche sur $H(J)$, et on a un iso. can. de $D$-multiplicités
\[
(155)\qquad D_J\simeq\overset{\deleted{\unsure{\sim}}}{D^!_d}\mathbin{\wedge}_{(\mathfrak S_d)_{D_d}}H(J).
\]

Soient \deleted{$J$} $(J,\underline{\unsure{c|d}})$ et $(J',\underline{c}')$ deux ens. ordonnés finis munis de fonctions codim, et
\[
(156)\qquad J'\mathrel{\overset{f}{\hookrightarrow}}J\deleted{\unsure{'}}
\]
un plong\textsuperscript{t} d'ens ordonnés, compatible aux fonctions codim, on trouve alors un \deleted{application} morphisme canonique
\[
(157)\qquad D_{J\deleted{\unsure{'}}}\longrightarrow D_{J'},
\]
je me demande \deleted{\unsure{sous forme}}\inserted{en termes} de quelles données suppl. \inserted{(si besoin est)} sur les $(D_d,\overset{\deleted{\unsure{\sim}}}{D^!_d})$ on peut reconstituer cette dépendance fonctorielle de $D_J$ \unsure{p.r.|p.v.} à $J$, via (155). \deleted{Considérons}\inserted{Comme la} \deleted{d'abord le plus petit élément $e_{J'}$ de $e$} dépendance fonctorielle en $J$ pour les \uline{iso}
\margin{109}
